\documentclass[a4paper, 11pt]{article}
\usepackage[utf8]{inputenc}
\usepackage[T1]{fontenc}
\usepackage[italian]{babel}
\usepackage{geometry}
\usepackage{graphicx}
\usepackage{tabularx}
\usepackage{fancyhdr}
\usepackage{hyperref}

% Margini
\geometry{top=3cm, bottom=3cm, left=2.5cm, right=2.5cm}

% Intestazione e piè di pagina
\pagestyle{fancy}
\fancyhf{}
\fancyhead[L]{\small \textit{0x06tudio} -- Verbale di Riunione}
\fancyhead[R]{\small Anno Accademico 2026/2027}
\fancyfoot[C]{\thepage}
\renewcommand{\headrulewidth}{0.4pt}

\begin{document}
\thispagestyle{plain}

\newcommand{\issue}[1]{\href{https://github.com/0x06tudio/docs/issues/#1}{\textit{#1}}}

\begin{minipage}{0.6\textwidth}
    \vspace{-2cm}
    {\huge \textbf{0x06tudio}} \\[0.5em]
    {\large Progetto didattico di Ingegneria del Software} \\[0.3em]
    {\normalsize Università degli Studi di Padova}
\end{minipage}
\begin{minipage}{0.4\textwidth}
    \raggedleft
    \vspace{0pt}
    \includegraphics[width=6cm]{../../media/logo.png}
    \vspace{1cm}
    \small
    \href{mailto:0x06tudio@gmail.com}{0x06tudio@gmail.com}
\end{minipage}

\vspace{0.5cm}

\begin{center}
    {\LARGE \textbf{Verbale Riunione}}
\end{center}

\vspace{0.5cm}


\begin{tabularx}{\textwidth}{l X}
    \textbf{Versione:} & 0.2.0 \\
    \textbf{Data:} & 2026/10/26 \\
    \textbf{Uso:} & Interno \\
    \textbf{Partecipanti (Presenti):} & Nicolò Beltrame, Pietro Brocadello, Alberto Cavalletto, Catalin Chirilov, Davide Colabove, Andrea Maggio \\
    \textbf{Destinatari:} & 0x06tudio \\
    \textbf{Modalità / Luogo:} & Online / Discord \\
    \textbf{Ora inizio:} & 15:30 \\
    \textbf{Ora fine:} & 17:15 \\
    \textbf{Durata totale:} & 1.45 ore \\
\end{tabularx}
\section*{Ordine del giorno}
\begin{enumerate}
    \item Formalizzazione del way-of-working e norme di gruppo;\vspace{-0.2cm}
    \item Prima assegnazione di ruoli e definizione della loro rotazione;\vspace{-0.2cm}
    \item Checklist delle attività in to-do;\vspace{-0.2cm}
    \item Assegnazione attività ai componenti;\vspace{-0.2cm}
    \item Recap scadenze ed inserimento a calendario di gruppo;\vspace{-0.2cm}
\end{enumerate}
\vspace{-0.5cm}

\section*{Registro delle modifiche}
\begin{tabularx}{\textwidth}{|c|c|X|c|c|}
    \hline
    \textbf{Versione} & \textbf{Data} & \textbf{Descrizione} & \textbf{Verificatore} & \textbf{Autore} \\ \hline
    0.2.0 & 2026/11/03 & Correzioni per nuova candidatura &  & Andrea Maggio \\ \hline
    0.1.1 & 2026/10/27 & Correzione numerazione ID & Nicolò Beltrame & Andrea Maggio \\ \hline
    0.1.0 & 2026/10/27 & Creazione e stesura bozza iniziale & Nicolò Beltrame & Davide Colabove \\ \hline
\end{tabularx}

\newpage
\tableofcontents
\newpage


\section{Resoconto della discussione}

\subsection{Punto 1: Formalizzazione del way-of-working e norme di gruppo}
Il team ha stabilito le norme base per il lavoro di gruppo. La frequenza degli incontri sarà di una volta a settimana fino all'inizio ufficiale del progetto, per poi intensificarsi a due incontri settimanali (con possibilità di ulteriore aumento in caso di necessità). Si è inoltre discussa l'idea di sviluppare un "sito vetrina" dedicato, con lo scopo di ospitare e rendere facilmente consultabile tutta la documentazione di progetto prodotta.

\subsection{Punto 2: Prima assegnazione di ruoli e definizione della loro rotazione}
Sono stati definiti i ruoli per il primo sprint di lavoro (S01):
\begin{itemize}
    \item \textbf{Responsabile:} Davide Colabove
    \item \textbf{Amministratore:} Catalin Chirilov
    \item \textbf{Analista:} Andrea Maggio
    \item \textbf{Progettista:} Pietro Brocadello
    \item \textbf{Programmatore:} Alberto Cavalletto
    \item \textbf{Verificatore:} Nicolò Beltrame
\end{itemize}
Il gruppo ha concordato che la durata dello sprint sarà indicativamente di due settimane. Al termine di ogni sprint avverrà la rotazione dei ruoli, che seguirà un meccanismo lineare di "shift +1" rispetto all'ordine attuale.

\subsection{Punto 3: Checklist delle attività in to-do}
Sono state discusse le attività operative necessarie per il completamento della candidatura e la gestione del gruppo. Le aree individuate comprendono: preparazione dei documenti di candidatura (Lettera di presentazione, Valutazione Capitolati, Preventivo Costi ed eventuale Studio di fattibilità), stesura dei verbali per i nuovi incontri interni ed esterni, preparazione dei materiali per il Diario di Bordo 2 (DdB 2) e gestione delle revisioni e delle firme.

\subsection{Punto 4: Assegnazione attività ai componenti}
I task individuati nel punto precedente sono stati assegnati ai membri del gruppo (il dettaglio è riportato nella tabella riepilogativa a fine verbale). Le poche attività ancora prive di un assegnatario esplicito verranno prese in carico internamente nei prossimi giorni.

\subsection{Punto 5: Recap scadenze ed inserimento a calendario di gruppo}
Al termine dell'incontro, tutte le scadenze relative ai task assegnati sono state ricontrollate e inserite nel calendario condiviso, al fine di garantire il corretto avanzamento dei lavori in vista della consegna della candidatura.


\section{Riepilogo delle decisioni e degli impegni (To-Do)}

\begin{tabularx}{\textwidth}{|c|X|X|}
    \hline
    \textbf{ID Decisione} & \textbf{Descrizione} & \textbf{Motivazione} \\ \hline
    D 3.1 & Frequenza incontri stabilita a 1 volta/settimana pre-progetto e 2 volte/settimana post-inizio. & Stabilire una cadenza regolare per gli allineamenti del gruppo\\ \hline
    D 3.2 & Rotazione dei ruoli fissata ogni 2 settimane tramite meccanismo di shift +1. & Requisito di progetto, tempo corrispondente ad uno sprint \\ \hline
\end{tabularx}

\vspace{2cm}

\noindent
\begin{tabularx}{\textwidth}{|c|X|X|c|X|}
    \hline
    \textbf{Task ID} & \textbf{Descrizione} & \textbf{Assegnato a} & \textbf{Scadenza} & \textbf{Issue} \\ \hline
    TD 3.1 & Stesura Lettera di presentazione & Pietro Brocadello & 2026/10/30 & \issue{7} \\ \hline
    TD 3.2 & Valutazione Capitolati (modifiche e completamento) & Catalin Chirilov & 2026/10/30 & \issue{22} \\ \hline
    TD 3.3 & Stesura Preventivo Costi & Alberto Cavalletto & 2026/10/30 & \issue{10} \\ \hline
    TD 3.4 & Stesura Verbale esterno (VarGroup) & Andrea Maggio & 2026/10/30 & \issue{23} \\ \hline
    TD 3.5 & Preparazione slide per Diario di Bordo 2 (DdB 2) & Davide Colabove & 2026/11/02 & \issue{14} \\ \hline
    TD 3.6 & Revisione documenti interni ed esterni & Nicolò Beltrame & 2026/10/30 & \issue{1}, \issue{2}, \issue{3}, \issue{4}, \issue{5}, \issue{6}, \issue{8}, \issue{9}, \issue{11}, \issue{12}, \issue{13}, \issue{15} \\ \hline
    TD 3.7 & Richiesta firme per verbali esterni & Nicolò Beltrame & 2026/10/30 & \issue{24} \\ \hline
\end{tabularx}

\end{document}